\section{Третья идея: Амортизированная нейросетевая дистилляция в O(1)}

\subsection{Мотивация: сжатие графа в единую нейросеть}

Хотя графовый планировщик Дейкстры показал высокую надежность, ему присущи два принципиальных недостатка: квадратичные затраты памяти $O(N^2)$ при масштабировании на гигантские лабиринты и вычислительная задержка ($\sim 1.2\,\text{мс}$ на шаг). 

Возникла исследовательская гипотеза: \textit{можно ли сжать глобальные рассуждения графового алгоритма в легковесную нейросеть-студент, которая за один прямой проход ($O(1)$) будет отображать пару $(s_t, B(g)) \to \hat{z}_{\text{cmd}}$ напрямую в низкоуровневую политику $\pi^\ell$?} Идейно такая модель принимает те же входы, что и базовый Single-Intention контроллер, но обучена не наивному спрямлению пути, а траекторному поведению учителя Дейкстры.

\subsection{Дифференцируемый функционал потерь в JAX}

Для обучения сети был собран датасет из 50~000 пар $(s_t, B(g)) \to z^*_{\text{teacher}}$, сгенерированных учителем. Чтобы сеть не просто копировала векторы (Behavior Cloning), но и учитывала физику робота и геометрию FB-пространства, я разработал 4-компонентный дифференцируемый лосс:
\addequation{
    \mathcal{L}(\theta) = \mathcal{L}_{\text{BC}} + 0.5 \mathcal{L}_{\text{Action}} - 0.02 \mathcal{L}_{\text{Reach}} - 0.05 \mathcal{L}_{\text{Goal}},
}{eq:distillation_loss}
где компоненты выполняют следующие физические функции:
\begin{itemize}[leftmargin=*,noitemsep]
    \item $\mathcal{L}_{\text{BC}} = (1 - \cos(\hat{z}, z^*)) + \text{MSE}(\hat{z}, z^*)$ — имитация направления вектора учителя.
    \item $\mathcal{L}_{\text{Action}} = \|\pi^\ell(s_t, \hat{z}) - a^*_{\text{teacher}}\|^2$ — прямое согласование действий в пространстве моторов (продифференцировано сквозь замороженную $\pi^\ell$).
    \item $\mathcal{L}_{\text{Reach}} = F(s_t, \hat{z})^\top \hat{z}$ — штраф за выбор недостижимых латентных направлений.
    \item $\mathcal{L}_{\text{Goal}} = \hat{z}^\top B(g)$ — глобальное притяжение к терминальной цели.
\end{itemize}

\begin{figure*}[t]
    \centering
    %!TEX root = ../main.tex
\begin{tikzpicture}[>=stealth, font=\small, scale=0.88, every node/.style={transform shape}]
    % Входные данные
    \node[draw=blue!80, fill=blue!10, rounded corners=4pt, minimum width=2.2cm, minimum height=0.9cm, align=center] (obs) at (0, 2.2) {Состояние $s_t$\\$\mathbb{R}^{29}$};
    \node[draw=orange!80, fill=orange!10, rounded corners=4pt, minimum width=2.2cm, minimum height=0.9cm, align=center] (goal) at (0, -0.2) {Цель $B(g)$\\$\mathbb{R}^{128}$};

    % Энкодеры
    \node[draw=blue!80, fill=blue!20, rounded corners=4pt, minimum width=2.4cm, minimum height=0.9cm, align=center] (enc_s) at (3.2, 2.2) {State Encoder\\MLP + SwiGLU};
    \node[draw=orange!80, fill=orange!20, rounded corners=4pt, minimum width=2.4cm, minimum height=0.9cm, align=center] (enc_g) at (3.2, -0.2) {Goal Encoder\\MLP + SwiGLU};

    % Блок внимания и гейта
    \node[draw=purple!80, fill=purple!15, rounded corners=6pt, minimum width=3.2cm, minimum height=1.6cm, align=center] (gate) at (7.0, 1.0) {\textbf{Gated Cross-Attention}\\$Q = E_s, \;\; K, V = E_g$\\$H = \text{Attn} \odot \sigma(W_g [E_s, E_g])$};

    % Проекция на сферу
    \node[draw=green!70!black, fill=green!15, rounded corners=4pt, minimum width=2.5cm, minimum height=0.9cm, align=center] (proj) at (10.8, 1.0) {\textbf{Проекция на $\mathbb{S}^{d-1}$}\\$\hat{z} = \sqrt{d} \cdot \frac{z}{\|z\|}$};

    % Выход
    \node[draw=red!70!black, fill=red!15, rounded corners=4pt, minimum width=2.2cm, minimum height=0.9cm, align=center] (out) at (13.8, 1.0) {\textbf{Интенция}\\$\hat{z}_{\text{cmd}} \to \pi^\ell$};

    \draw[->, thick] (obs) -- (enc_s);
    \draw[->, thick] (goal) -- (enc_g);
    \draw[->, thick] (enc_s) -- (gate);
    \draw[->, thick] (enc_g) -- (gate);
    \draw[->, thick] (gate) -- (proj);
    \draw[->, thick] (proj) -- (out);
\end{tikzpicture}

    \caption{Архитектура амортизированного планировщика Distilled JAX Gated-Attention с дифференцируемым лоссом в $O(1)$.}
    \label{fig:distilled_jax}
\end{figure*}

\subsection{Исследование нейросетевых архитектур}

В рамках дистилляции было исследовано 4 архитектурных решения: Standard MLP, ResNet-ECA (1D-канальное внимание), Dense-ECA и Gated Cross-Attention со SwiGLU.

\begin{table}[H]
    \centering
    \footnotesize
    \caption{Сравнение архитектур дистилляции на Medium}
    \label{tab:arch_comparison}
    \begin{tabular}{lcccc}
        \toprule
        \textbf{Архитектура} & \textbf{Парам.} & \textbf{Cos-Sim} & \textbf{SR (\%)} & \textbf{Время} \\
        \midrule
        Standard MLP & 312K & 0.891 & 78.0\% & 0.12 мс \\
        Dense-ECA & 458K & 0.914 & 82.0\% & 0.16 мс \\
        ResNet-ECA & 524K & 0.908 & 80.0\% & 0.15 мс \\
        \textbf{Gated Cross-Attn} & \textbf{680K} & \textbf{0.942} & \textbf{86.0\%} & \textbf{0.18 мс} \\
        \bottomrule
    \end{tabular}
\end{table}

Архитектура Gated Cross-Attention показала лучший результат, достигнув **86.0\%** успешности (превзойдя учителя Дейкстры на +8\% за счет сглаживания дискретных артефактов графа) при инференсе всего **0.18 мс** на шаг.
